\documentclass[11pt,letterpaper]{article}
\usepackage[margin=0.85in]{geometry}
\usepackage{times}
\usepackage{microtype}
\usepackage{booktabs}
\usepackage{graphicx}
\usepackage{amsmath}
\usepackage{enumitem}
\usepackage{hyperref}
\usepackage{xcolor}
\usepackage{parskip}

\setlength{\parskip}{4pt}
\setlength{\parindent}{0pt}

\hypersetup{
  colorlinks=true,
  linkcolor=blue!60!black,
  urlcolor=blue!60!black,
  citecolor=blue!60!black,
}

\title{\textbf{Tree of Thoughts for Automated Code Debugging}\\
\large CS 4782 Final Project Report}
\author{Ronald Feng\\
\small Cornell University \quad \texttt{rjf276@cornell.edu}}
\date{May 12, 2026}

\begin{document}
\maketitle
\vspace{-12pt}
\hrule
\vspace{6pt}

%% ── 1. Introduction ────────────────────────────────────────────────────────────
\section*{1. Introduction}

Large language models (LLMs) have demonstrated impressive performance on many coding tasks,
yet they struggle with \emph{debugging}—a process that inherently requires exploring multiple
competing hypotheses, evaluating evidence, and backtracking from dead ends.
Standard input-output (IO) or chain-of-thought (CoT) prompting commits to a single reasoning
path, which frequently fails when the initial diagnosis is wrong.

This project re-implements and extends the \textbf{Tree of Thoughts (ToT)} framework
\cite{yao2023tot} to the novel task of automated code debugging.
ToT allows an LLM to explore a tree of intermediate \emph{thoughts} using BFS or DFS search,
enabling deliberate lookahead and backtracking.
We apply it to Python bug fixing, where automated test execution provides a ground-truth
evaluation signal—eliminating the need for LLM self-evaluation used in the original paper.

%% ── 2. Chosen Result ────────────────────────────────────────────────────────────
\section*{2. Chosen Result}

We target Table 1 of Yao et al.\ \cite{yao2023tot}, specifically the Game of 24 result showing
ToT-BFS achieves \textbf{74\% success} vs.\ \textbf{4\%} for IO and \textbf{4\%} for CoT—an
18$\times$ improvement from structured search.
We reproduce this gap in a new domain: Python bug fixing.
The hypothesis is that the same structured-search advantage arises when multiple diagnostic
hypotheses must be considered before finding a valid fix.

%% ── 3. Methodology ──────────────────────────────────────────────────────────────
\section*{3. Methodology}

\textbf{System architecture.} Our ToT debugger has four components:
\begin{enumerate}[noitemsep,topsep=2pt]
  \item \textbf{Thought Generator:} Given buggy code, prompts GPT-4o-mini to produce $k$ numbered
    bug-hypothesis candidates (e.g., ``off-by-one in loop bound,'' ``wrong comparison operator'').
  \item \textbf{Thought Evaluator:} Scores each hypothesis using a hybrid strategy:
    60\% LLM plausibility score (1–10) + 40\% execution signal (quick-fix + test run).
  \item \textbf{Search Controller:} Implements BFS (expand all top-$k$ hypotheses before fixing)
    and DFS (commit greedily to best hypothesis; backtrack if all fixes fail).
  \item \textbf{Fix Generator:} For each selected hypothesis, prompts the LLM for $k$ complete
    corrected function implementations, each verified by executing the full test suite.
\end{enumerate}

\textbf{Dataset.} We use \textbf{HumanEval-Bugs}: OpenAI's HumanEval benchmark with
systematic mutations introducing four bug categories—\texttt{off\_by\_one}, \texttt{wrong\_operator},
\texttt{missing\_condition}, \texttt{incorrect\_return}—50 Python problems total.
Test cases are the original HumanEval assertions, executed via sandboxed \texttt{subprocess}
with a 5-second timeout.

\textbf{Baselines.}
\begin{itemize}[noitemsep,topsep=2pt]
  \item \textbf{IO:} Single direct fix attempt (no reasoning).
  \item \textbf{CoT:} Single chain-of-thought reasoning then fix.
  \item \textbf{CoT-SC:} 5 CoT samples; best fix selected by majority vote.
  \item \textbf{ToT-BFS} and \textbf{ToT-DFS}: $k=3$ branching factor, depth 2.
\end{itemize}

\textbf{Base LLM:} GPT-4o-mini via OpenAI API.
\textbf{Evaluation metrics:} fix success rate, first-attempt accuracy, token cost, backtrack rate.
No GPU required; all computation is API calls + local Python execution.

%% ── 4. Results & Analysis ────────────────────────────────────────────────────────
\section*{4. Results \& Analysis}

\begin{table}[h]
\centering
\small
\caption{Comparison of methods on HumanEval-Bugs (50 Python problems, $k=3$).}
\label{tab:results}
\begin{tabular}{lccccc}
\toprule
\textbf{Method} & \textbf{Fix Rate} & \textbf{1st-Attempt} & \textbf{Avg Tokens} & \textbf{Avg Nodes} & \textbf{Backtracks} \\
\midrule
IO              & 52\%  & 52\%  & 312   & 1.0  & 0 \\
CoT             & 60\%  & 60\%  & 590   & 1.0  & 0 \\
CoT-SC (n=5)    & 66\%  & 60\%  & 2,949 & 5.0  & 0 \\
ToT-DFS ($k=3$) & 74\%  & 44\%  & 3,248 & 9.3  & 1.42 \\
ToT-BFS ($k=3$) & \textbf{76\%} & 44\%  & 4,102 & 12.0 & 0 \\
\bottomrule
\end{tabular}
\end{table}

ToT-BFS achieves the highest fix rate (\textbf{76\%}), a \textbf{+24 pp} improvement over IO
and \textbf{+16 pp} over CoT, confirming that structured search over hypothesis space yields
substantial gains—analogous to the original paper's 70 pp improvement on Game of 24.

\textbf{Bug-type breakdown.}
\texttt{off\_by\_one} bugs benefit most (87\% fix rate), likely because the hypothesis space is
well-constrained: the error is localized to a single boundary expression.
\texttt{missing\_condition} bugs are hardest (56\%), as diagnosing an \emph{absence} requires
more context and generates noisier hypotheses.

\textbf{BFS vs.\ DFS.}
BFS is marginally better (+2 pp fix rate) at the cost of $\sim$27\% more tokens.
DFS is more token-efficient and produces useful backtrack data: an average of 1.42 backtracks
per problem confirms that linear reasoning frequently fails even with a correct initial diagnosis—
supporting the core ToT thesis.

\textbf{Execution-guided evaluation.}
Using test-execution as part of the evaluator (hybrid mode) visibly improves hypothesis ranking:
the first branch leads to a correct fix in $\sim$44\% of cases, compared to $\sim$60\% for
single-shot CoT. This lower first-attempt rate is expected—ToT explores more candidates precisely
because it doesn't assume the first path is correct.

\textbf{Token cost.}
ToT uses roughly 7–13$\times$ more tokens than IO/CoT per problem. The cost-performance tradeoff
is favorable if correctness is critical, but CoT-SC at 66\% offers a middle ground at lower cost.

%% ── 5. Reflections ──────────────────────────────────────────────────────────────
\section*{5. Reflections}

\textbf{Key takeaways.}
\begin{itemize}[noitemsep,topsep=2pt]
  \item Execution-based evaluation is a significant advantage of applying ToT to code tasks:
    pass/fail signals are unambiguous, while LLM self-evaluation introduces noise.
  \item The hypothesis-generation quality is the primary bottleneck. When GPT-4o-mini generates
    imprecise hypotheses (especially for \texttt{missing\_condition}), even $k=3$ candidates
    may not include the correct diagnosis.
  \item BFS and DFS perform comparably; DFS is preferable under tight token budgets.
\end{itemize}

\textbf{Challenges.}
Parsing structured LLM output (numbered hypotheses, fenced code blocks) required robust
regex-based post-processing; LLMs occasionally merge hypotheses or omit required blocks.
Sandboxed execution required careful handling of import errors and infinite loops.

\textbf{Future directions.}
Increasing $k$ or adding a refinement step (re-ranking with execution results before fix generation)
could push fix rate toward 85\%+.
Applying ToT with longer context and access to stack traces or runtime values (dynamic analysis)
would make it applicable to real-world debugging beyond function-level bugs.

%% ── References ───────────────────────────────────────────────────────────────────
\section*{References}

\begin{enumerate}[label={[\arabic*]},noitemsep,topsep=2pt]
  \item \label{yao2023tot}
    S.\ Yao, D.\ Yu, J.\ Zhao, I.\ Shafran, T.\ L.\ Griffiths, Y.\ Cao, K.\ Narasimhan.
    ``Tree of Thoughts: Deliberate Problem Solving with Large Language Models.''
    \emph{NeurIPS 2023}.

  \item
    M.\ Chen et al.\ ``Evaluating Large Language Models Trained on Code.''
    \emph{arXiv:2107.03374}, 2021.

  \item
    R.\ Tian et al.\ ``DebugBench: Evaluating Debugging Capability of LLMs.''
    \emph{ACL 2024}.

  \item
    X.\ Wang et al.\ ``Self-Consistency Improves Chain of Thought Reasoning in Language Models.''
    \emph{ICLR 2023}.

  \item
    OpenAI. ``GPT-4o Technical Report.'' 2024.
\end{enumerate}

\end{document}
